\section{张量场}\label{sec:tensor-fields}

向量场的定义可以推广到一般类型的张量场. 为此, 将 $M$ 中一点 $p$ 处的切空间 $T_pM$ 的对偶空间记为 $T_p^*M$ (通常称为 $M$ 在 $p$ 处的\textbf{余切空间}).

\begin{RiemannDefinition}\label{def:tensor-field}
一个 \textbf{$(k,m)$-张量场}是给每个点 $p\in M$ 指派一个张量 $T\in\mathcal{T}^{k,m}(T_p^*M,T_pM)$ 的映射.
\end{RiemannDefinition}

\begin{RiemannExample}\label{ex:vector-field-as-a-tensor-field}
向量场是一个 $(0,1)$-张量场 (或 1 重逆变张量场), 即给每个点 $p\in M$ 指派一个 1 重逆变张量 $X_p\in T_pM$ 的映射.
\end{RiemannExample}

\begin{RiemannExample}\label{ex:differential-is-covariant-tensor-field}
设 $f:M\to\mathbb{R}$ 是可微函数. 可以定义一个 $(1,0)$-张量场 $\mathrm{d}f$, 它将每个点 $p\in M$ 映到 $(\mathrm{d}f)_p$, 其中
\[
  (\mathrm{d}f)_p:T_pM\to\mathbb{R}
\]
是 $f$ 在 $p$ 处的导数. 这个张量场称为 $f$ 的\textbf{微分}. 对任意 $v\in T_pM$, 有 $(\mathrm{d}f)_p(v)=v\cdot f$ (即 $f$ 在 $p$ 处沿向量 $v$ 的方向导数). 考虑坐标系 $x:W\to\mathbb{R}^n$, 可以写成 $v=\sum_{i=1}^nv^i\bigl(\frac{\partial}{\partial x^i}\bigr)_p$, 所以
\[
  (\mathrm{d}f)_p(v)=\sum_i v^i\frac{\partial\widehat f}{\partial x^i}(x(p)),
\]
其中 $\widehat f=f\circ x^{-1}$. 取坐标函数 $x^i:W\to\mathbb{R}$, 可以得到定义在 $W$ 上的 1-形式 $\mathrm{d}x^i$. 它们满足
\[
  (\mathrm{d}x^i)_p\left(\left(\frac{\partial}{\partial x^j}\right)_p\right)=\delta_{ij}
\]
因而构成每个余切空间 $T_p^*M$ 的一组基, 与 $T_pM$ 的坐标基 $\bigl\{\bigl(\frac{\partial}{\partial x^1}\bigr)_p,\ldots,\bigl(\frac{\partial}{\partial x^n}\bigr)_p\bigr\}$ 对偶. 所以 $W$ 上任意 $(1,0)$-张量场可写为 $\omega=\sum_i\omega_i\,\mathrm{d}x^i$, 其中 $\omega_i:W\to\mathbb{R}$ 满足 $\omega_i(p)=\omega_p\bigl(\bigl(\frac{\partial}{\partial x^i}\bigr)_p\bigr)$. 特别地, $\mathrm{d}f$ 可以写成通常的形式
\[
  (\mathrm{d}f)_p=\sum_{i=1}^n\frac{\partial\widehat f}{\partial x^i}(x(p))(\mathrm{d}x^i)_p.
\]
\end{RiemannExample}

\begin{RiemannRemark}\label{rem:cotangent-bundle}
与切丛的构造类似, 可以考虑所有余切空间的不交并, 得到流形
\[
  T^*M=\bigcup_{p\in M}T_p^*M
\]
称为 $M$ 的\textbf{余切丛}. 注意, 一个 $(1,0)$-张量场就是由
\[
  p\mapsto\omega_p\in T_p^*M
\]
定义的从 $M$ 到 $T^*M$ 的映射. 这个构造很容易推广到任意类型的张量场.
\end{RiemannRemark}

显然, $(k,m)$-张量场的空间是向量空间, 因为 $(k,m)$-张量的线性组合仍然是 $(k,m)$-张量. 若 $W$ 是 $M$ 的一个坐标邻域, 则已知 $\{(\mathrm{d}x^i)_p\}$ 是 $T_p^*M$ 的一组基, 而 $\bigl\{\bigl(\frac{\partial}{\partial x^i}\bigr)_p\bigr\}$ 是 $T_pM$ 的一组基. 因此, 一个 $(k,m)$-张量场 $T$ 在点 $p\in W$ 处的值可以写成张量
\begin{align*}
  T_p&=\sum a_{i_1\ldots i_k}^{j_1\cdots j_m}(p)
  (\mathrm{d}x^{i_1})_p\otimes\cdots\otimes(\mathrm{d}x^{i_k})_p\\
  &\qquad\otimes\left(\frac{\partial}{\partial x^{j_1}}\right)_p
  \otimes\cdots\otimes\left(\frac{\partial}{\partial x^{j_m}}\right)_p
\end{align*}
其中 $a_{i_1\ldots i_k}^{j_1\ldots j_m}:W\to\mathbb{R}$ 这些函数在每个点 $p\in W$ 处给出 $T_p$ 相对于 $T_p^*M$ 和 $T_pM$ 的这些基的分量. 如同向量场的情形, 若对极大图册的所有坐标系, 这些函数都可微, 就称该张量场\textbf{可微}. 同样, 只需考虑一个图册的坐标系, 因为所有重叠映射均可微 [参见 Exercise~\ref{ex:tensor-fields}(\ref{ex:components-of-tensor-under-coordinate-change})].

\begin{RiemannExample}\label{ex:smooth-function-has-smooth-differential}
光滑函数 $f:M\to\mathbb{R}$ 的微分显然是一个可微 $(1,0)$-张量场, 因为在给定坐标系 $x:W\to\mathbb{R}^n$ 下, 它的分量 $\frac{\partial\widehat f}{\partial x^i}\circ x$ 都是光滑的.
\end{RiemannExample}

协变张量的一个重要运算是沿光滑映射的\textbf{拉回}.

\begin{RiemannDefinition}\label{def:pullback-of-covariant-tensor-field}
设 $f:M\to N$ 是光滑流形之间的可微映射. 则 $N$ 上的每个可微 $k$ 重协变张量场 $T$ 都按以下方式在 $M$ 上定义一个 $k$ 重协变张量场 $f^*T$:
\[
  (f^*T)_p(v_1,\ldots,v_k)=T_{f(p)}((\mathrm{d}f)_pv_1,\ldots,(\mathrm{d}f)_pv_k),
\]
其中 $v_1,\ldots,v_k\in T_pM$.
\end{RiemannDefinition}

\begin{RiemannRemark}\label{rem:pointwise-and-algebraic-properties-of-tensor-pullback}
注意, $(f^*T)_p$ 正是 $T_{f(p)}$ 在线性映射 $(\mathrm{d}f)_p^*:\mathcal{T}^k(T_{f(p)}^*N)\to\mathcal{T}^k(T_p^*M)$ 下的像, 这一映射由 $(\mathrm{d}f)_p:T_pM\to T_{f(p)}N$ 诱导 (参见第~\ref{sec:tensors}~节). 因此, 对所有 $\alpha,\beta\in\mathbb{R}$ 及所有适当的协变张量场 $T,S$, 有 $f^*(\alpha T+\beta S)=\alpha(f^*T)+\beta(f^*S)$ 和 $f^*(T\otimes S)=(f^*T)\otimes(f^*S)$. 我们将在 Exercise~\ref{ex:tensor-fields}(\ref{ex:smooth-pullback-of-covariant-tensor}) 中看到, 可微协变张量场的拉回仍是可微协变张量场.
\end{RiemannRemark}

\begin{RiemannExercise}\label{ex:tensor-fields}
\begin{enumerate}[label=(\arabic*),ref=\arabic*,leftmargin=*]
\item\label{ex:components-of-tensor-under-coordinate-change} 求张量场在两个重叠坐标系中的坐标函数之间的关系.
\item\label{ex:smooth-pullback-of-covariant-tensor} 证明可微协变张量场的拉回仍是可微协变张量场.
\item\label{ex:lie-derivative-of-tensors} (\emph{张量场的 Lie 导数}) 给定向量场 $X\in\mathfrak{X}(M)$, 定义\textbf{$k$ 重协变张量场 $T$ 沿 $X$ 的 Lie 导数}为
\[
  L_XT:=\left.\frac{\mathrm{d}}{\mathrm{d}t}(\psi_t^*T)\right|_{t=0},
\]
其中 $\psi_t=F(\cdot,t)$, $F$ 是 $X$ 在 $p$ 处的局部流.
\begin{enumerate}[label=(\alph*),leftmargin=*]
\item 证明
\begin{align*}
  L_X(T(Y_1,\ldots,Y_k))
  &=(L_XT)(Y_1,\ldots,Y_k)\\
  &\quad+T(L_XY_1,\ldots,Y_k)+\cdots+T(Y_1,\ldots,L_XY_k),
\end{align*}
即证明
\begin{align*}
  X\cdot(T(Y_1,\ldots,Y_k))
  &=(L_XT)(Y_1,\ldots,Y_k)\\
  &\quad+T([X,Y_1],\ldots,Y_k)+\cdots+T(Y_1,\ldots,[X,Y_k]),
\end{align*}
对所有向量场 $Y_1,\ldots,Y_k$ 成立 [参见第~\ref{chap:differentiable-manifolds}~章的 Exercise~\ref{ex:vector-fields}(\ref{ex:lie-derivative-of-function}) 和 Exercise~\ref{ex:vector-fields}(\ref{ex:lie-derivative-of-vector-field})].
\item 你会如何定义 $(k,m)$-张量场的 Lie 导数?
\end{enumerate}
\end{enumerate}
\end{RiemannExercise}
